\section{Conclusiones}

En términos generales, el desarrollo de este proyecto representó un reto integral, abarcando desde la concepción de una arquitectura desde cero hasta la toma de decisiones críticas sobre las herramientas tecnológicas a utilizar. Como primer gran proyecto de la materia, supuso un excelente acercamiento al ciclo de vida del desarrollo de software con aplicaciones del mundo real. Aunque el producto final se encuentra en una etapa temprana, la experiencia resultó sumamente desafiante y enriquecedora en el ámbito académico y pre-profesional.

Un aprendizaje fundamental durante este proceso fue la gestión del tiempo y la capacidad de autoestudio. En el desarrollo de software es común enfrentarse a una curva de aprendizaje pronunciada, donde el exceso de investigación puede generar un estancamiento práctico. La rápida evolución de las herramientas exige que los estudiantes y profesionales de la computación desarrollen la habilidad de aprender sobre la marcha y adaptarse continuamente; de lo contrario, el conocimiento técnico queda obsoleto con rapidez.

Asimismo, este proyecto destacó la vital importancia del trabajo colaborativo. Enfrentar problemas técnicos complejos de forma colectiva, apoyándose en el conocimiento y la experiencia de los compañeros de clase, demostró ser una herramienta invaluable. En una industria que en ocasiones fomenta el aislamiento y la individualidad, desarrollar redes de apoyo y un sentido de comunidad facilita enormemente la resolución de problemas lógicos y fortalece las habilidades interpersonales.

Finalmente, el sistema Cliente-Servidor construido posee una base sólida con un alto potencial de escalabilidad. Como trabajo futuro, se podrían proponer diversas expansiones inspiradas en aplicaciones de mensajería modernas, tales como:

\begin{itemize}
    \item Implementación de fotos de perfil vinculadas a cada usuario.
    \item Soporte para el envío de archivos multimedia, incluyendo imágenes, notas de voz o fragmentos de video.
\end{itemize}

Para lograr estas implementaciones, se reconoce que sería indispensable realizar modificaciones sustanciales al protocolo de comunicación, adaptando la serialización JSON y la lógica de ambos extremos (cliente y servidor) para que sean capaces de procesar, transmitir y renderizar este nuevo volumen de datos estructurados.